\section{Coding with a decreasing dependency bound}
\label{sec:bootstrap}

We prove Theorem~\ref{thm:monoids}.  The two monoids in that theorem
must contain information about one another: the finitely presented monoid
will multiply representatives of the decidable monoid, while the latter
will contain a coded copy of the former.  We arrange this using a program
whose evaluations on words of length $n$ depend only on its values at
strictly smaller lengths.  The other ingredient is a family of separated
words that supplies the finite-tuple selectors.

\subsection{The bounded compiler}

Fix a finite alphabet $\Sigma$ containing the disjoint triples
\[
 (\mathtt b_T,\mathtt a_T,\mathtt c_T),\qquad
 (\mathtt b_G,\mathtt a_G,\mathtt c_G)
\]
and the six further letters
$\mathtt z,\mathtt d,\mathtt l,\mathtt r,\mathtt h,\zeta$.
Fix a total order with $\mathtt a_T<\mathtt c_T$ and
$\mathtt a_G<\mathtt c_G$.  Words are ordered by \emph{shortlex}:
first by length and then lexicographically.  We use an acceptable
programming system for partial functions $\Sigma^*\to\Sigma^*$, writing
$f_\eta$ for the function with index $\eta$.  In particular, universal
evaluation and effective substitution of parameters are available;
integers and tuples are represented by fixed computable word encodings.

The following finite construction is proved in
Appendix~\ref{sec:compiler}.  Its conditional statement is essential:
the program index determines the presentation without running the
program. The evaluator is also fixed independently of a length bound;
its termination and compatibility with equations will be proved from
the program's behavior within that bound.

\begin{lemma}[Bounded compiler]\label{lem:compiler}
From a program index $\eta$ one can compute a finite monoid presentation
$K_\eta$ with a zero letter $\Omega$, a finite presentation alphabet
$X_\eta$, and an additive grade $\gamma$ on words, such that
$0\le\gamma(x)\le M=2$ for every letter and $\gamma(\Omega)=0$.
There is a partial word algorithm $N_\eta$, independent of any budget,
with the following properties.

The alphabet contains letters $\alpha,\mathsf Q,\omega$ and right
anchor letters $a_R$ for $a\in\Sigma$, all mutually distinct and
different from $\Omega$. For a word $u$, let $u_R$ replace each
letter $a$ by $a_R$, and set
\[
 \operatorname{raw}(u)=\alpha\mathsf Q u_R\omega\qquad(u\in\Sigma^*).
\]
Unconditionally, the algorithm fixes the empty word, $\Omega$, and
every literal word $\operatorname{raw}(u)$, and sends every word
containing $\Omega$ to $\Omega$. On an input word $a$, its only calls
to $f_\eta$ are on words of length at most $\gamma(a)$; apart from
those calls it performs finite computations.

Suppose that, for some integer $m\ge0$,
\begin{align}
 &f_\eta(w)\text{ is defined and }|f_\eta(w)|\le |w|
       &&(|w|\le m),\label{eq:budget-total}\\
 &f_\eta(xf_\eta(w)y)=f_\eta(xwy)
       &&(|xwy|\le m).\label{eq:budget-context}
\end{align}
Then $N_\eta$ terminates on all words of grade at most $m$.
Each such word is joined to its literal output by a finite derivation
in the presentation whose intermediate grades do not exceed the input
grade.  A single presentation equation, used in any context, preserves
the output whenever both endpoints have grade at most $m$.
Consequently
\begin{equation}\label{eq:bounded-context}
 N_\eta(u)=N_\eta(v)\quad\Longrightarrow\quad
 N_\eta(xuy)=N_\eta(xvy)
 \quad\text{if }\gamma(xuy),\gamma(xvy)\le m.
\end{equation}

If
\eqref{eq:budget-total}--\eqref{eq:budget-context} hold for all $m$,
then $N_\eta$ decides equality in $K_\eta$ and
\begin{equation}\label{eq:raw-product}
 N_\eta\bigl(\operatorname{raw}(u)\operatorname{raw}(v)\bigr)=\operatorname{raw}\bigl(f_\eta(uv)\bigr).
\end{equation}
The presentation is computed without evaluating $f_\eta$.
\end{lemma}

For clarity, \eqref{eq:bounded-context} follows from the preceding
derivation statement.  Transplant the paths from $u$ and $v$ to their
common output into $x(-)y$.  Neither contextual path leaves the budget,
so one-step invariance identifies their outputs.  A general machine
encoding without this bounded conclusion would not suffice below.
The raw singleton rule has a different role: it will distinguish the
elements used for the injection $H\to K$. Multiplication of two raw
words, rather than evaluation of a singleton, will encode the product
in $H$ through \eqref{eq:raw-product}.

\subsection{Codes and finite rule tables}

First consider any index $\eta$.  The auxiliary monoid $\mathcal P$
has decidable equality by Lemma~\ref{lem:partial-affine}.

Let $Y$ be the disjoint union of $X_\eta$ and a fixed finite generating
alphabet for $\mathcal P$.  For a $Y$-word $v$, write $v_K,v_{\mathcal P}$
for its projections, obtained by deleting the letters of the other
alphabet.  Its \emph{tested class} is zero if $N_\eta(v_K)=\Omega$;
otherwise it is the pair consisting of $N_\eta(v_K)$ and the value of
$v_{\mathcal P}$ in $\mathcal P$.  This is at present a partial test.
For the second type of code, fix a finite monoid generating alphabet
$Z$ for $\widehat G$, including the inverses of a finite group
generating set, and test words by their values in $\widehat G$.
This second test is total by the hypothesis of Theorem~\ref{thm:monoids}.

Choose computably an integer $D_0>M$ such that
$2^{D_0-1}\ge\max\{|Y|,|Z|\}$.  Encode the letters of $Y$, in their
chosen order, by distinct words beginning with $\mathtt b_T$ and
followed by $D_0-1$ letters from
$\{\mathtt a_T,\mathtt c_T\}$, in lexicographic order.  Encode $Z$
in the same way using the other triple.  Write $C(v)$ for concatenation
of the appropriate codes and put $C(\varepsilon)=\varepsilon$.
The type of a nonempty code sequence is unambiguous.  Codes preserve
shortlex comparison and have length $D_0$ per generator.

Enumerate all nonempty finite ordered tuples
\[
 (w_1^{(j)},\ldots,w_{s_j}^{(j)})\qquad(j\ge1)
\]
of distinct $\Sigma$-words avoiding $\zeta$; the empty word is allowed
as a member.  Choose integers $N_j$ recursively so that
\begin{equation}\label{eq:envelope-spacing}
 N_j>j,\qquad N_j>\max_i|w_i^{(j)}|,\qquad
 N_j>2N_h+2+\max_i|w_i^{(h)}|\quad(h<j).
\end{equation}
These inequalities concern \emph{all potential} earlier envelopes,
regardless of which rules will be selected.  The increasing sequence
$N_j$ is computable without performing any tested-class evaluation.

We now define a partial procedure $F(\eta,w)$.  Put $n=|w|$ and form
the following finite rule table $\mathcal R_n$.
\begin{enumerate}
\item Include $\zeta a\to\zeta$ and $a\zeta\to\zeta$ for every
letter $a\in\Sigma$.  For each nonempty $Y$-word $v$ with
$D_0|v|\le n$, test its class.  If it is zero, include
$C(v)\to\zeta$.  Otherwise find the shortlex least word $v'$ of the
same tested class among all words of length at most $|v|$, including
the empty word; include $C(v)\to C(v')$ if $v'\ne v$.
Do the same for $Z$-words using equality in $\widehat G$, with no
zero case.  Call these the \emph{static rules}.

\item Process the tuples with $N_j\le n$ in order.  If every member
of tuple $j$ has no reducible occurrence under the static rules and
previously selected rules, include
\begin{equation}\label{eq:selector-rules}
 \mathtt d\mathtt l^{N_j}w_i^{(j)}\mathtt r^{N_j}\mathtt h
 \longrightarrow
 \begin{cases}
 \varepsilon,&i=1,\\
 \zeta,&i>1.
 \end{cases}
\end{equation}
Some of these left sides may be longer than $n$; they simply cannot
act on the input.

\item Reduce $w$ until no rule applies, using a fixed deterministic
choice of occurrence and rule.
\end{enumerate}
Whenever the tests terminate, the last step terminates: every rule
decreases shortlex.  In particular a nonempty code reduced to $\zeta$
shortens, since $D_0>2$.  All other non-code rules shorten as well.
Thus $F$ is uniformly partial computable.

\subsection{The fixed point and its length induction}

Effective substitution of parameters gives a total computable index
transformation $\tau(\eta)$ for the function $F(\eta,-)$.  Kleene's
Recursion Theorem~\cite{Kleene38} supplies an index with
\begin{equation}\label{eq:fixed-point}
 f_\eta(w)=F(\eta,w)
\end{equation}
as partial functions.  Here is the short parameterization argument.
Let $U$ be universal evaluation and let $d(a)$ be an effectively
obtained index for the partial function
$w\mapsto U(U(a,a),w)$.  Choose an index $b$ computing the total
function $a\mapsto\tau(d(a))$.  For $\eta=d(b)$,
\[
 U(\eta,w)=U(U(b,b),w)=U(\tau(d(b)),w)=F(\eta,w).
\]
Only now fix the presentation alphabet and code length belonging to
this particular $\eta$.  There is no advance bound on the machine's
size or working space.  Write $f=f_\eta$ and $K=K_\eta$.

\begin{lemma}\label{lem:bootstrap}
The function $f$ is total and length-nonincreasing.  The rule tables
are stable on bounded-length words: for $k\le n$, $\mathcal R_k$
and $\mathcal R_n$ have the same rules whose left sides have length
at most $k$.  Each $\mathcal R_n$ is confluent on words of length at
most $n$, and
\begin{equation}\label{eq:f-context}
 f(xf(w)y)=f(xwy)\qquad(x,w,y\in\Sigma^*).
\end{equation}
\end{lemma}

\begin{proof}
We prove all bounded assertions together by induction on $n$.
For $n=0$, there are no tested classes or tuple decisions, and the
empty word is unchanged.  Suppose $n>0$ and the assertions are known
at smaller bounds.  Every use of $N_\eta$ in a static test, including
each candidate in a least-representative search, uses a $Y$-word of length at most
$\lfloor n/D_0\rfloor$.  Its $K$-projection has grade at most
\begin{equation}\label{eq:strict-dependency}
 m=M\lfloor n/D_0\rfloor<n.
\end{equation}
The induction hypothesis supplies
\eqref{eq:budget-total}--\eqref{eq:budget-context} at $m$.
Lemma~\ref{lem:compiler} therefore makes all these tests terminate.
There are only finitely many tests and tuple decisions, and the
subsequent shortlex-decreasing reduction terminates.

We first establish stability.  For a fixed generator word $v$, its
static rule is determined by the same tests and the same candidates
of length at most $|v|$ at every larger bound.  Consider a tuple with
$N_j\le k\le n$.  Each tested body has length less than $N_j$, hence
less than $k$.  New static left sides longer than $k$ cannot occur
in it.  Induction over earlier tuple numbers shows that the previously
selected rules that can occur in the bodies also agree.  Hence the
decision for tuple $j$ agrees at both bounds.  A new tuple with
$N_j>k$ has every selector side longer than $k$, so cannot add a rule
on a word of length at most $k$.  This proves stability, including
the decisions for selected rules too long to act at the smaller bound.

It remains to prove local confluence for $\mathcal R_n$.  All
reductions stay inside the length-at-most-$n$ ball.  Disjoint
occurrences can be reduced in either order. If the word contains
$\zeta$, every branch still contains it: absorption retains $\zeta$,
and no other rule has a left side containing that letter. Absorption joins all
branches at $\zeta$.  We check the other overlaps explicitly.

\smallskip\noindent
\emph{Selectors from different stages.}
Let the older left side have total length $L$, and let the newer one
be $\mathtt d\mathtt l^N w\mathtt r^N\mathtt h$.
By \eqref{eq:envelope-spacing}, $L<N$ and $|w|<N$.
A nonempty old suffix cannot equal a new prefix: the former ends in
$\mathtt h$, while a prefix of length at most $L$ lies in the initial
$\mathtt d\mathtt l^N$.  An old prefix cannot equal a new suffix,
since their first letters are respectively $\mathtt d$ and either
$\mathtt r$ or $\mathtt h$.  These exclude both orientations of a
proper boundary overlap.  For containment, an old occurrence cannot
start at the new initial $\mathtt d$, because its terminal
$\mathtt h$ would lie in the initial $\mathtt l$-run.  Any other
start must be in the body.  It cannot leave the body: it cannot end
in the final $\mathtt r$-run, and reaching the final $\mathtt h$
would require crossing a run longer than the entire old side.
Thus it would lie wholly in the body, contrary to the newer tuple's
irreducibility test.

\smallskip\noindent
\emph{Selectors from one stage.}
A second start strictly inside one side must lie in its body.
There are then fewer than $N$ body letters left before the final
$\mathtt r$-run, so that start cannot be followed by $N$ consecutive
$\mathtt l$'s.  If two sides start at the same place and their body
lengths differ, the terminal $\mathtt h$ of the shorter falls in the
final $\mathtt r$-run of the longer, because both body lengths are
less than $N$.  Equal body lengths force equal bodies, hence the same
rule, as tuple members are distinct.  This also covers an empty body.

\smallskip\noindent
\emph{A code and a selector.}
Code letters are disjoint from the envelope markers.  An overlapping
code side would therefore have to lie wholly inside the selector's
body.  The static irreducibility test excluded this.  All static sides
short enough to occur there were already present when the tuple was
processed, because the body's length is less than $N_j$.

\smallskip\noindent
\emph{Two code sides.}
They have the same type and aligned starts: within a code sequence,
$\mathtt b_T$ or $\mathtt b_G$ occurs only at a block boundary.
Their union is therefore an entire code sequence $C(v)$, of length
at most $n$.  For group codes, a replacement preserves the value of
this union by group multiplication.  For product codes, replacing
a subword by a nonzero tested-equal word preserves the tested class
of the union by \eqref{eq:bounded-context} and equality in
$\mathcal P$.  Indeed each old or new union has generator length
at most $\lfloor n/D_0\rfloor$, and hence its $K$-projection has
grade at most $m$.  We use this bound separately for both endpoints:
a shorter generator word need not have smaller actual $K$-grade.

If a subword tests zero, apply \eqref{eq:bounded-context} to its
$K$-projection and $\Omega$, whose grade is zero.  The contextual
word containing $\Omega$ has output $\Omega$, so the union also
tests zero.  Hence if one branch introduces $\zeta$, it reduces the
whole union to $\zeta$, and the other either does the same or retains
a zero-class code sequence, to which its whole-word zero rule applies.
That remaining sequence cannot be empty, since
$N_\eta(\varepsilon)=\varepsilon\ne\Omega$.
The same argument covers a zero-class union when both first
replacements are code sequences.

If the union class is nonzero, let $v_*$ be its least representative
among words of length at most $|v|$.  Either branch yields a word
$v_b$ of the same class with $|v_b|\le|v|$.  Since $v_b$ is one of
the candidates defining $v_*$, we have
$|v_*|\le|v_b|$, and $v_*$ is also the least candidate within this
smaller bound.  Each branch therefore reduces by its whole-word rule
to $C(v_*)$, unless it already equals that word.  The representative
$v_*$ may be empty; no empty left side is needed.  These joining
reductions can be made on the union before acting on any newly
adjacent outside letters.

This exhausts the local peaks.  Termination and local confluence give
confluence on the ball by Newman's Lemma~\cite{Newman42}.
Indeed, by well-founded induction, assume every proper descendant of
a word has a unique normal form.  Any two first steps have a common
descendant by local confluence; the induction hypothesis makes their
normal forms equal, proving uniqueness for the original word and
hence confluence.  Finally, if $|xwy|\le n$,
perform inside $x(-)y$ the normalization of $w$ computed at bound
$|w|$.  By stability every rule used belongs to $\mathcal R_n$, and
length never increases.  Confluence gives
$f(xf(w)y)=f(xwy)$.  This establishes the budget hypotheses at $n$
and completes the induction.
\end{proof}

\subsection{The two monoids and their embeddings}

Let $H$ be the monoid on $\Sigma$ whose equations are the union of
all the rule tables.  Every word is related to its computed normal
form, which is irreducible for this union by stability.
Conversely, a contextual use of any one
equation preserves that form: choose a bound containing both
endpoints and the equation, and apply stability and confluence from
Lemma~\ref{lem:bootstrap}.  Thus two words represent the same element
of $H$ if and only if their literal $f$-outputs agree.  This proves
decidability.  The empty word and the single letter $\zeta$ are
distinct irreducibles; they represent the identity and the absorbing
zero of $H$.

The budget hypotheses now hold at every length, so
Lemma~\ref{lem:compiler} gives an equality test for $K$.  In
particular its identity and zero are distinct.  The tested product
classes are now the actual elements of
\[
 T=\bigl((K\setminus\{\Omega\})\times\mathcal P\bigr)\cup\{0_T\},
\]
with coordinatewise multiplication and every pair having first
coordinate $\Omega$ collapsed to zero.
Here a letter of $X_\eta$ acts in the first coordinate and a
generator of $\mathcal P$ in the second; the letter $\Omega$ maps to
$0_T$.  Thus $Y$ generates $T$.

For a nonzero element of $T$, take its least shortlex $Y$-representative
$v$.  The code $C(v)$ is irreducible: a subcode rule would either
force the element to be zero or produce a smaller representative;
no selector or absorption side occurs in a pure code word.
Every other representative reduces by its whole-word rule to this
same code, and every zero-class code reduces to $\zeta$.
Consequently coding defines an injective zero-preserving monoid map
\[
 \kappa:T\longrightarrow H.
\]
It is unital because the least representative of the identity is
empty, and it is multiplicative because codes concatenate.
The identical argument for $Z$-words embeds $\widehat G$ unitally in
$H$.  Its image consists of units, since the products of mutually
inverse group words have the empty representative.  This includes
the trivial group and does not require a nonempty generating set.

Let $z\in H$ be represented by the separator $\mathtt z$.
For least representatives $v,v'$ of nonzero elements of $T$, the words
\begin{equation}\label{eq:separator-words}
 C(v)\mathtt z C(v')
\end{equation}
are irreducible.  No code crosses the separator; each code portion
is already irreducible; and there are no selector markers or
$\zeta$.  These literal words recover the ordered pair, including
when one or both representatives are empty.  They are nonzero
because none equals $\zeta$.  This proves the separator assertion
of Theorem~\ref{thm:monoids}.

\subsection{Selectors and multiplication by raw words}

Take any ordered tuple of distinct nonzero elements
$h_1,\ldots,h_s$ of $H$, and let $w_1,\ldots,w_s$ be their normal
words.  These words avoid $\zeta$, since any word containing it
reduces to the zero word.  Their tuple occurs in the enumeration,
say at stage $j$.  A word irreducible for the final rule set is
irreducible for the static and selected rules present at stage $j$.
Hence this tuple passes its test.  With
\[
 L=\mathtt d\mathtt l^{N_j},\qquad
 R=\mathtt r^{N_j}\mathtt h
\]
as elements of $H$, the selected equations give
$Lh_1R=1$ and $Lh_iR=\zeta$ for $i>1$.  The argument includes a tuple
containing the identity, whose normal representative is empty.

Finally define a multiplicative map in the opposite direction by
\[
 c:H\longrightarrow K,\qquad c([w])=[\operatorname{raw}(f(w))].
\]
The normal-form characterization of $H$ makes this well-defined.
Every raw singleton is its own literal $N_\eta$-output.  Distinct
normal words therefore give distinct elements of $K$, and none of
these elements is its identity or its zero.  For normal words $u,v$,
\eqref{eq:raw-product} gives
\[
 c([u])c([v])=[\operatorname{raw}(f(uv))]=c([uv]).
\]
Thus $c$ is injective and multiplicative, as required.  It is not
unital: even $c(1)=[\operatorname{raw}(\varepsilon)]$ is distinct from $1_K$.
Also $c(\zeta)\ne\Omega$; it is absorbing only within $c(H)$.
The finite presentation of $K$ came from the compiler, while the
finite alphabet and normal-form procedure give the claimed finite
generation and decidability of $H$.  All conclusions of
Theorem~\ref{thm:monoids} are proved.
